\section{B. Good reduction and inertia characters}
\subsection{Explicit counterexamples to the mod-$p$ good-reduction criterion}

\begin{theorem}[B-RD1: a counterexample for every odd prime]\label{thm:B-RD1}
\status{Proved.}
For every prime $p\geq3$, there is a complete discretely valued field $K$
of characteristic zero, residue field $\overline{\mathbf F}_p$, and
absolute ramification index
\[
 e(K)=\frac{p^2-1}{2},
\]
containing $\mu_p$, and an elliptic curve $E/K$ whose full $p$-torsion
is $K$-rational and whose reduction is additive, potentially good and
supersingular. Its semistability defect is exactly $2$. Its inertia
character is $-2$ on the nontrivial element of the finite inertia image.
\end{theorem}
\begin{proof}
We give the field and the curve by an explicit algebraic construction.
Put $k=\overline{\mathbf F}_p$, $R=W(k)$ and $F=\operatorname{Frac}R$.
Let $m=(p-1)/2$ and choose any root $\lambda\in k$ of
\[
 H_p(T)=\sum_{i=0}^{m}\binom mi^2T^i.
\]
Such a root exists since $m\geq1$. It is neither $0$ nor $1$:
$H_p(0)=1$ and
$H_p(1)=\binom{2m}{m}\equiv(-1)^m\pmod p$.
The coefficient calculation for Frobenius on the one-dimensional space
$H^1(\mathcal O)$ of a Legendre elliptic curve gives $(-1)^mH_p(\lambda)$
as its Hasse--Witt scalar. Thus
\[
 E_0/R:\quad y^2=x(x-1)(x-[\lambda])
\]
is smooth and has supersingular special fibre. Here $[\lambda]$ is the
Teichmuller representative. The discriminant is
$16[\lambda]^2(1-[\lambda])^2$, an $R$-unit.

Let $L=F(E_0[p])$. We now verify its ramification without invoking a
Breuil--Kisin dictionary. Every nonzero geometric $p$-torsion point
specializes to the identity, since the special fibre is supersingular.
For a local parameter $T$ of the formal group at that identity, its
height is two. Consequently $[p](T)/T$ has constant coefficient $p$,
its first coefficient which is a unit modulo $p$ is that of
$T^{p^2-1}$, and all lower nonconstant coefficients are divisible by $p$.
Weierstrass preparation gives a polynomial of degree $p^2-1$ which is
Eisenstein over $R$. Thus every nonzero $P\in E_0[p]$ generates a
completely ramified extension $F(P)/F$ of degree $p^2-1$. This degree is
prime to $p$, so the extension is tame. Since $F$ has algebraically
closed residue field, its unique tame extension of that degree is
Galois cyclic and contains every such point. Hence
\[
 \operatorname{Gal}(L/F)\simeq C_{p^2-1}.
\]
Its action on $E_0[p]\setminus\{0\}$ is free and transitive. The unique
involution in this cyclic group acts as $-1$: an involution in
$\mathrm{GL}_2(\mathbf F_p)$ with any eigenvalue $+1$ would fix a
nonzero vector.

Let $H$ be this order-two subgroup and put $K=L^H$. The Weil pairing
shows $\mu_p\subset L$, and $H$ acts on it by the determinant of $-1$,
which is $1$. Thus $\mu_p\subset K$. We have $e(K)=(p^2-1)/2$.
Choose any nonzero $P\in E_0[p]$, and put $\delta=y(P)^2\in K^\times$.
A nonzero odd-order point has $y(P)\ne0$, so $L=K(\sqrt\delta)$.
Write $E_0$ as $y^2=x^3+a_2x^2+a_4x+a_6$ and define
\[
 E:\quad y^2=x^3+\delta a_2x^2+\delta^2a_4x+\delta^3a_6.
\]
Over $L$, the map $(x,y)\mapsto(\delta x,\delta\sqrt\delta\,y)$
identifies $E_0$ with $E$. For every $Q\in E_0[p]$, both $x(Q)$ and
$\sqrt\delta\,y(Q)$ are fixed by $H$. Hence every point of $E[p]$
is $K$-rational. Equivalently, its representation is the tensor product
of the scalar quadratic representation on $E_0[p]$ and its inverse.

The curve $E_0$ has good reduction over $K$, so for every prime
$\ell\ne p$ inertia acts trivially on $T_\ell(E_0)$. On $T_\ell(E)$ it
acts by the nontrivial ramified quadratic character of $L/K$. The
Neron--Ogg--Shafarevich criterion\cite[Theorem~1]{SerreTate1968Good} therefore excludes good reduction.
It has good supersingular reduction over $L$, so its bad reduction is
additive and its minimal good-reduction inertia quotient has order $2$.
The trace on its covariant Dieudonne module is $-2$ at the involution,
since that involution acts as $[-1]$ on the special elliptic curve.
\end{proof}

\begin{theorem}[B-RD2: full $3$-torsion already over $\mathbf Q_3(\zeta_3)$]\label{thm:B-RD2}
\status{Proved, with an exact computation.}
Let $K=\mathbf Q_3(\zeta)$, where $\zeta^2+\zeta+1=0$. Then $K$ has
absolute ramification index $2$ and residue field $\mathbf F_3$. The curve
\[
 E/K:\quad y^2+y=x^3
\]
has $E[3]\subset E(K)$, additive potentially supersingular reduction,
and semistability defect exactly $6$. Its inertia character takes the
value $-2$ on the element of order $2$ of its finite inertia image.
\end{theorem}
\begin{proof}
The discriminant and $j$-invariant are $\Delta=-27$ and $j=0$.
The third division polynomial is
\[
 \psi_3(x)=3x(x^3+1).
\]
Besides the identity, the following are the eight distinct $3$-torsion
points:
\[
 (0,0),\ (0,-1),\qquad
 (x,y)\quad\text{with }x\in\{-1,-\zeta,-\zeta^2\},\quad
 y\in\{\zeta,\zeta^2\}.
\]
They exhaust $E[3]$, which has degree nine in characteristic zero.
The exact program \texttt{verify\_p3\_full\_torsion.py} checks all these
points, their orders, all $81$ sums, and that $(0,0),(-1,\zeta)$ form a
basis. Its recorded output is part of the computation certificate.

The displayed equation has $v_K(\Delta)=6$. Discriminants of two
Weierstrass equations of the same curve differ by a twelfth power.
Thus no equation over $K$ can have discriminant valuation zero, and
$E$ does not have good reduction.

We construct good reduction after ramification degree six. Make first
the unramified quadratic extension $K'=K(i)$, with $i^2=-1$, and choose
$\sqrt3\in K'$. Adjoin $r$ with $r^3=-1/4$, and then $u$ with
$u^2=r\sqrt3$. The polynomial for $s=r+1$ is
\[
 s^3-3s^2+3s-3/4.
\]
Its Newton polygon has the single slope $-2/3$ for the valuation
normalized on $K'$, so $K'(r)/K'$ is completely ramified of degree three.
The element $r\sqrt3$ has odd valuation in that field, so adjoining $u$
gives a completely ramified quadratic extension. Put $L=K'(r,u)$.
After completing the square and setting
\[
 x=r+u^2X,\qquad y+\tfrac12=u^3Y,
\]
the equation becomes
\[
 Y^2=X^3+\frac{3r}{u^2}X^2+X.
\]
The middle coefficient has positive valuation and the special fibre is
$Y^2=X^3+X$, which is smooth and supersingular in characteristic three.
Thus $E$ has good reduction over an extension of ramification degree six.

We verify minimality, including possible unramified extensions. Let
$F/K$ have good reduction. The discriminant valuation implies
$2\mid e(F/K)$. In a good integral Weierstrass equation over $F$, complete
the square; this is integral since $2$ is a unit. The resulting equation
has the form $Y^2=X^3+b_2X^2+b_4X+b_6$. Comparison with
$y^2=x^3+1/4$ uses $x=r_0+u_0^2X$, $y=u_0^3Y$, and forces
$v_3(u_0)=1/4$. Integrality of $b_6$ gives
\[
 v_3(r_0^3+1/4)\geq3/2.
\]
It follows that $r_0$ is a unit with residue $-1$. Writing $s_0=r_0+1$,
if $v_3(s_0)<1/3$ the term $s_0^3$ uniquely has lowest valuation in
$s_0^3-3s_0^2+3s_0-3/4$; if $v_3(s_0)>1/3$ the constant term uniquely
has lowest valuation. Both alternatives give valuation at most one.
Hence $v_3(s_0)=1/3$, so $3\mid e(F/\mathbf Q_3)$, and therefore
$3\mid e(F/K)$. We conclude $6\mid e(F/K)$, proving the defect is six.

Over the unramified extension $K'$, the extension just constructed is
Galois: $r\mapsto\zeta r$, $u\mapsto\zeta^2u$ has order three and
commutes with $u\mapsto-u$. Its inertia group is cyclic of order six.
The latter involution induces $Y\mapsto-Y$ on the good special fibre,
that is $[-1]$, so its trace is $-2$.
\end{proof}

\begin{proposition}[B-RD3: the ordinary case of R3(a)]\label{prop:B-RD3}
\status{Proved.}
Let $K$ be complete discretely valued of characteristic zero, with
perfect residue field of characteristic $p\geq3$, and assume
$\mu_p\subset K$. If an abelian variety $A/K$ has potentially good
ordinary reduction and $A[p]$ is unramified, then $A$ has good reduction.
There is no restriction on $\dim A$.
\end{proposition}
\begin{proof}
We may pass to the completed maximal unramified extension: good
reduction is equivalent before and after this passage by the
Neron--Ogg--Shafarevich criterion. The residue field is now algebraically
closed. Choose a finite Galois extension $L/K$ giving good reduction,
and let $\mathcal A/\mathcal O_L$ be the resulting abelian scheme.
For its ordinary special fibre, the connected subgroup
$\mathcal A[p]^0$ is of multiplicative type and its etale quotient has
rank $p^g$. Indeed the Cartier dual of the connected subgroup has
etale special fibre, hence is etale by Nakayama's lemma applied to its
relative differentials. Over the strictly henselian ring
$\mathcal O_L$, it is constant; thus
$\mathcal A[p]^0\simeq\mu_p^g$.

Every inertia element acts semilinearly on this subgroup by its action
on the base, composed with a matrix in $\mathrm{GL}_g(\mathbf F_p)$.
The base action fixes $\mu_p$, and its generic action on $A[p]$ is
trivial. The matrix is therefore the identity. The action on the etale
quotient is trivial for the same reason. Over the algebraically closed
residue field the connected--etale sequence splits, and there are no
nonzero homomorphisms between the connected and etale factors in either
direction. Thus inertia acts identically on $\mathcal A_k[p]$.

If $\alpha$ is one of the resulting finite-order automorphisms of
$\mathcal A_k$, then $\alpha-1$ factors through $[p]$, so
$\alpha=1+p\beta$ for some $\beta\in\operatorname{End}(\mathcal A_k)$.
The kernel of reduction modulo $p$ in the units of this endomorphism
ring has no nontrivial finite-order elements for $p\geq3$. To see this,
pass to its finite free $\mathbf Z_p$-completion. If
$\alpha=1+p^a\gamma$ with $a\geq1$ and $\gamma\notin p\operatorname{End}$,
the binomial expansion shows that a prime-to-$p$ power has the same
leading valuation, and the $p$-th power has leading valuation $a+1$;
neither can be one. This proves the assertion for any finite order.
Hence inertia is trivial on the entire special abelian variety.
Proper smooth base change now makes $T_\ell(A)$ unramified for
$\ell\ne p$, and the same good-reduction criterion concludes.
\end{proof}

\begin{corollary}[B-RD6: the manuscript's cas test is negative]\label{cor:B-RD6}
\status{Refuted, for every odd prime.}
For every $p\geq3$ there is a field $K$ as in Theorem~\ref{thm:B-RD1}
and two elliptic curves $E,E'/K$ with potentially good reduction and
isomorphic $p$-torsion, for which the inertia traces differ on an element
of order prime to $p$.
\end{corollary}
\begin{proof}
Use $E$ and $K$ from Theorem~\ref{thm:B-RD1}. Choose an ordinary
elliptic curve $C/\overline{\mathbf F}_p$; for example choose a Legendre
parameter outside the finite root set of $H_p$ and outside $\{0,1\}$.
The ordinary Serre--Tate coordinates of
\cite[Theorem~2.1]{Katz1981SerreTate} give its canonical deformation
over $W(\overline{\mathbf F}_p)$, whose $p$-divisible group is
\[
 \mu_{p^\infty}\oplus\mathbf Q_p/\mathbf Z_p.
\]
This formal elliptic scheme algebraizes: its identity section provides
the compatible relatively ample bundles $\mathcal O(3[0])$, so the
formal existence theorem applies. Let $E'$ be its generic fibre after
extension to $K$. Then $E'$ has good reduction, and $\mu_p\subset K$
gives $E'[p]\simeq(\mathbf Z/p)^2\simeq E[p]$ over $K$.
The good curve has inertia trace $2$; the bad quadratic twist has trace
$-2$ at the inertia involution. That involution has order prime to $p$.
\end{proof}

\subsection{The first torsion level: a sharp tame ramification range}

\begin{theorem}[B-RD7: the first-level tame range]\label{thm:B-RD7}
\status{Proved; independently audited by the parent agent.}
Let $p\geq5$, and permit arbitrary perfect residue fields of characteristic
$p$. There exists a complete discretely valued $K$ of characteristic zero
and absolute ramification index $e$, and an elliptic curve $E/K$ with
$E[p]\subset E(K)$, additive potentially good reduction, and semistability
defect $m$, if and only if
\[
 m\in\{2,3,4,6\},\qquad m\mid(p+1),\qquad
 (p-1)\mid e,\qquad e\geq\frac{p^2-1}{m}.
\]
In the existence direction one may take the residue field to be
$\overline{\mathbf F}_p$. In particular the set of possible $e$ is
\[
 \left\{e\in(p-1)\mathbf Z_{>0}:
 e\geq\frac{p^2-1}{m_{\max}(p)}\right\},
\]
where $m_{\max}(p)=2,6,4,6$ for $p\equiv1,5,7,11\pmod {12}$,
respectively. For $p=3$, the same assertion with a specified tame defect
$m=2$ or $4$ holds: defect $2$ occurs exactly at even $e\geq4$, and
defect $4$ occurs at every positive even $e$. No assertion about wild
3-primary inertia is included in the latter statement.
\end{theorem}

\begin{proof}
We first prove necessity. Completion of the maximal unramified extension
does not change $e$ or the inertia image. The ordinary case is excluded
by Proposition~\ref{prop:B-RD3}. For $p\geq5$, potential-good inertia is
tame by \cite[Corollary~2 of Theorem~2]{SerreTate1968Good}. Thus the
minimal good-reduction extension $L/K$ is cyclic of order $m$, with
supersingular reduction $C$. Over an algebraically closed field of
characteristic greater than three, writing $C$ in short Weierstrass form
shows that its origin-preserving automorphisms have the form
$(x,y)\mapsto(u^2x,u^3y)$, with $u^4a=a$ and $u^6b=b$.
The possibilities are $m=2,3,4,6$. The curves with $j=0$ are supersingular
exactly when $p\equiv-1\pmod3$, and those with $j=1728$ exactly when
$p\equiv-1\pmod4$. Both assertions follow directly by extracting the
coefficient of $x^{p-1}$ from $(x^3+b)^{(p-1)/2}$ and
$(x^3+ax)^{(p-1)/2}$, respectively. Consequently $m\mid p+1$.
The same discussion for tame inertia in characteristic three gives
$m=2$ or $4$.

Normalize the valuation $v_L$ by $v_L(L^\times)=\mathbf Z$.
An average over the tame cyclic action provides a formal parameter $t$
at the identity of the good model which is an exact eigenvector.
If its eigenvalue is a primitive $m$th root $\zeta^q$, then every
$K$-rational nonzero $p$-torsion point $P$ satisfies
\[
 v_L(t(P))\equiv q\pmod m,\qquad \gcd(q,m)=1.
\]
All these valuations are positive, since $C$ is supersingular. Choose a
point $P$ whose parameter has maximum valuation $A$, and a point $Q$
whose parameter has minimum valuation $B$. If $A>B$ they are independent;
if $A=B$, choose any basis. The $p-1$ nonzero multiples of $P$ have
valuation $A$. The other $p(p-1)$ points have valuation $B$: if $A>B$,
this follows from the linear term of the formal group law; if $A=B$,
it follows from the assumed common maximum and minimum.
The constant term of $[p](t)/t$ is $p$, and its distinguished degree is
$p^2-1$. Weierstrass preparation therefore gives
\[
 m e=v_L(p)=(p-1)(A+pB).                         \tag{B.7.1}
\]
Here $A\geq B\geq1$ and $A\equiv B\equiv q\pmod m$.
Since $p\equiv-1\pmod m$, the displayed identity implies both
$(p-1)\mid e$ and $me\geq p^2-1$.

We give a deformation-theoretic converse with the equivariance made
explicit. Put $k=\overline{\mathbf F}_p$. Choose a supersingular elliptic
curve $C/k$ with an automorphism $\alpha$ of order $m$. For $m=3,6$
choose $j=0$, for $m=4$ choose $j=1728$ (the unique supersingular
$j$ in characteristic three is also zero), and for $m=2$ choose any
supersingular curve. The derivative of $\alpha$ is a primitive $m$th
root; a finite automorphism of order prime to $p$ with derivative one
is the identity, by averaging a local coordinate.

The framed deformation functor of $C$, with full Drinfeld $p$-level
structure, is represented by a complete regular local ring $R$ of
dimension two. If $P,Q$ are the universal labelled basis, their formal
coordinates form a regular system of parameters. These are the
precise statements of \cite[Theorem~2.1.2 and
Proposition~2.2.3(i)]{CarterStrauch2021LubinTate}; Serre--Tate identifies
the height-two formal-group deformation with the elliptic deformation
\cite[Theorem~1.2.1]{Katz1981SerreTate}. The universal formal elliptic
curve algebraizes: the identity gives the compatible ample line
bundles $\mathcal O(3[0])$, so the projective formal existence theorem
applies over the complete noetherian ring $R$.

Changing the framing by $\alpha$ gives an action of $C_m$ on this
functor, hence a semilinear action on the universal elliptic curve.
It preserves the labels $P,Q$. For any initial formal coordinate $t_0$,
replace it by
\[
 t=\frac1m\sum_{i=0}^{m-1}\zeta^{-i}\alpha^{i*}t_0,
\]
where $\zeta$ is the derivative eigenvalue. The linear coefficient is
a unit, so $t$ remains a coordinate and $\alpha^*t=\zeta t$.
Thus $x=t(P)$ and $y=t(Q)$ are regular parameters satisfying exactly
$\alpha(x)=\zeta x$, $\alpha(y)=\zeta y$. The action fixes $W(k)$.
Cohen's presentation theorem gives an equivariant surjection
\[
 W(k)[[X,Y]]\longrightarrow R,\qquad X\mapsto x,\quad Y\mapsto y,
\]
where $\alpha$ multiplies $X,Y$ by $\zeta$.

We detail its single relation, to avoid assuming that an arbitrary
monomial specialization has the intended ramification. The full-level
factorization gives, in $R$,
\[
 p=u\prod_{(a,b)\in\mathbf F_p^2\setminus\{0\}}
          t([a]P+[b]Q),\qquad u\in R^\times.
\]
Lift $u$ to a unit $U$ and each factor to a power series $T_{a,b}$
in $W(k)[[X,Y]]$ by using the multiplication and addition power
series of the formal group, lifting their coefficients. These lifts
have zero constant term and linear term $aX+bY$ (using integral
representatives for $a,b$); crucially, $T_{a,0}$ is divisible by $X$
with unit quotient when $a\ne0$. Then
\[
 F_0=p-U\prod_{(a,b)\ne0}T_{a,b}
\]
lies in the kernel and generates it: it has coefficient one on $p$
modulo the square of the maximal ideal of $W(k)[[X,Y]]$.
Replace $F_0$ by its $C_m$-average $F$. This still generates the kernel,
has the same linear term $p$, and is invariant.

Write $t=e/(p-1)$ and choose
\[
 B=1,\qquad A=1+m\left(t-\frac{p+1}{m}\right).
\]
These are positive integers, $A\geq B$, and $A\equiv B\equiv1\pmod m$.
Reduce the relation modulo $p$, and give $X,Y$ weights $A,B$. If $A>B$, the weighted initial term of the
product in $F_0$ is a nonzero scalar times
$X^{p-1}Y^{p(p-1)}$, of weight
$d=(p-1)(A+pB)$. If $A=B$, the weighted initial term is the nonzero
homogeneous polynomial $\prod_{(a,b)\ne0}(aX+bY)$ of degree $p^2-1$.
In either case that initial form is invariant, since its total degree
is $p^2-1$, divisible by $m$. Averaging therefore preserves it.
Choose a Teichmuller unit $c$ arbitrarily when $A>B$, and with
$\bar c\notin\mathbf F_p$ when $A=B$. Then
\[
 R_L=W(k)[[\pi]]/(F(\pi^A,c\pi^B))
\]
is a complete discrete valuation ring with uniformizer $\pi$ and
$v_L(p)=d$: modulo $p$ its defining equation has first nonzero term
of degree $d$, while its coefficient on $p$ is a unit. This also
proves directly that the specified specialization exists.

The action $\pi\mapsto\zeta\pi$ preserves this equation, by invariance
of $F$ and $A\equiv B\equiv1\pmod m$. It has order $m$, acts trivially
on the residue field, and makes the map $R\to R_L$ equivariant.
Let $L=\operatorname{Frac}R_L$ and $K=L^{C_m}$. Then $L/K$ is tame
of ramification degree $m$, and
\[
 e(K)=d/m=e.
\]
The generic elliptic curve descends to $K$. Since the labels of the
level structure were preserved, its full $p$-torsion is rational over
$K$. The inertia action on the good special fibre is the faithful
order-$m$ action $\alpha$. The Neron--Ogg--Shafarevich criterion gives
exactly semistability defect $m$, so reduction is additive. Finally the
Weil pairing implies $\mu_p\subset K$. This proves every sufficiency
assertion made in the theorem, including $p=3$ and $m=2,4$.
\end{proof}

\subsection{A missing wild semistability defect in characteristic three}

\begin{theorem}[B-RD8: defect twelve with full rational $3$-torsion]\label{thm:B-RD8}
\status{Proved; exact computation recorded.}
The list $2,3,4,6$ in B.2 is incomplete in residue characteristic three.
Put
\[
 K=\mathbf Q_3(i,\zeta),\qquad i^2=-1,\quad
 \zeta^2+\zeta+1=0,\qquad \pi=\zeta-1,\quad t=1+i\pi.
\]
Then $K$ has absolute ramification index $2$ and residue field
$\mathbf F_9$. The elliptic curve
\[
 E:\quad y^2=x^3-27t(t^3+8)x+54(t^6-20t^3-8)
\]
has full $3$-torsion over $K$ and additive potentially supersingular
reduction, with semistability defect exactly $12$.
\end{theorem}
\begin{proof}
Write $a=-27t(t^3+8)$ and $b=54(t^6-20t^3-8)$.
Set $\delta=1+2\zeta=3+2\pi$, so $\delta^2=-3$.
The eight nonzero $3$-torsion points are
\[
 \left(-9t^2,\ \pm12\delta(t^3-1)\right)
\]
and, for $\epsilon\in\{1,\zeta,\zeta^2\}$,
\[
 \left(3t^2+12\epsilon t+12\epsilon^2,
       \ \pm36(\epsilon t^2+\epsilon^2t+1)\right).
\]
Indeed their four $x$-coordinates give the factorization
\[
 3x^4+6ax^2+12bx-a^2
 =3(x+9t^2)\prod_{\epsilon^3=1}
       (x-3t^2-12\epsilon t-12\epsilon^2).
\]
Substitution verifies their $y$-coordinates; the points are distinct
and exhaust the eight nonzero geometric $3$-torsion points.
The exact certificate \texttt{verify\_p3\_defect12.py} performs these
calculations in
$\mathbf Q[i,\pi]/(i^2+1,\pi^2+3\pi+3)$, verifies the curve equation
and the tripling identity for every point, and verifies independence
of points with different $x$-coordinates. No numerical approximation
is used.

Normalize $v_K(\pi)=1$. Direct calculation gives
\[
 \Delta=2^{12}3^9(t^3-1)^3,\quad
 v_K(t^3-1)=v_K(t^3+8)=3,\quad v_K(\Delta)=27.
\]
In particular $E$ is not good over $K$, since the valuation of a
Weierstrass discriminant can change only by a multiple of $12$.
We next both construct good reduction and prove the required minimality.

In the $2$-division polynomial $x^3+ax+b$, substitute
$x=3\pi^2 Z$. Dividing by $(3\pi^2)^3$ gives
\[
 Z^3+A Z+B,
 \qquad A=4i\pi+3\pi+6,\quad B=-12i+20\pi+30.
\]
Both $A$ and $B$ have valuation one. The polynomial is Eisenstein,
so a root $z$ generates a totally ramified cubic extension $K_1/K$,
and $v_K(z)=1/3$. Set $x_0=3\pi^2 z$. Then
\[
 v_K(x_0)=13/3,\qquad v_K(3x_0^2+a)=9.
\]
The second equality can also be read from the derivative
$3z^2+A$, whose valuation is one. Adjoin $u$ satisfying
\[
 u^4=3x_0^2+a.
\]
In the normalized valuation of $K_1$, the right side has value $27$,
prime to four. Thus this is a totally ramified degree-four extension;
for $L=K_1(u)$ we have $e(L/K)=12$ and $v_K(u)=9/4$.
The substitution $x=x_0+u^2X$, $y=u^3Y$ gives
\[
 Y^2=X^3+\frac{3x_0}{u^2}X^2+X.
\]
The middle coefficient has valuation $11/6>0$. The special fibre is
$Y^2=X^3+X$, smooth and supersingular in characteristic three.
Therefore good reduction is obtained with ramification degree twelve.

Conversely, let $F/K$ be any finite extension over which $E$ has good
reduction. Its discriminant valuation implies
$12\mid27e(F/K)$, hence $4\mid e(F/K)$. Since $2$ is invertible,
the $2$-torsion of a good elliptic scheme is finite etale; all its
points are defined over the maximal unramified extension of $F$.
The Eisenstein $2$-division root has $K$-valuation $1/3$ after the
indicated scaling, so $3\mid e(F/K)$. Consequently
$12\mid e(F/K)$. Together with the constructed extension, this proves
that the semistability defect is exactly twelve.
\end{proof}

\subsection{The complete first-level range in residue characteristic three}

\begin{lemma}[B-RD9: a variable-ramification Hesse family]\label{lem:B-RD9}
\status{Proved.}
Let $K$ have algebraically closed residue field of characteristic three,
absolute ramification index $e$, and a uniformizer $\varpi$. Suppose
$\mu_3\subset K$. Choose an integer $r$ with
$0<2r<e$ and $3\nmid e+r$, and put $t=1+\varpi^r$. The Hesse
Weierstrass curve of Theorem~\ref{thm:B-RD8} has full rational
$3$-torsion and semistability defect
\[
 m=\frac{12}{\gcd(4,e+r)}.
\]
\end{lemma}
\begin{proof}
The torsion formulas in Theorem~\ref{thm:B-RD8} are identities in $t$,
so apply unchanged. Write $w=\varpi^r$ and first scale the $2$-division
polynomial by $x=3X$. After translating $X=w^2+Z$, the resulting cubic is
\[
 Z^3+3w^2Z^2-(12w^3+18w^2+36w+27)Z
       +12w^4-36w^3-117w^2-108w-54.
\]
Its coefficients, from constant through quadratic, have valuations
$e+4r,e+3r,e+2r$. The inequalities $2r<e$ make the displayed leading
terms unique. Its Newton polygon has the single slope
$-(e+4r)/3$. Because $3\nmid e+r$, adjoining a root $z$ gives a
completely ramified cubic extension $K_1/K$. Put $x_0=3(w^2+z)$.
Then
\[
 v_K(x_0)=e+2r,\qquad v_K(3x_0^2+a)=3(e+r).
\]
The derivative computation giving the second equality has a unique
lowest term, of valuation $e+3r$, before multiplying by $3^2$.

Choose $u$ with $u^4=3x_0^2+a$. Since the residue field is algebraically
closed and $4$ is invertible, adjoining this root has ramification
degree $4/\gcd(4,e+r)$ over $K_1$. The equation obtained from
$x=x_0+u^2X$, $y=u^3Y$ is
\[
 Y^2=X^3+(3x_0/u^2)X^2+X,
\]
and the middle coefficient has valuation $(e+r)/2>0$. Thus it is good.
Conversely $v_K(\Delta)=9(e+r)$, which forces the tame factor
$4/\gcd(4,e+r)$ in any good-reduction ramification degree. The
$2$-division cubic forces the additional factor three, exactly as in
Theorem~\ref{thm:B-RD8}. This proves minimality.
\end{proof}

\begin{theorem}[B-RD10: all first-level defects at $p=3$]\label{thm:B-RD10}
\status{Proved; independently audited by the parent agent.}
Allowing arbitrary perfect residue fields, a potentially good additive
elliptic curve with full rational $3$-torsion exists at every positive
even absolute ramification index. The precise possible defects are
\[
\begin{array}{c|c}
\text{defect}&\text{absolute ramification indices}\\\hline
2& e\in2\mathbf Z,\ e\geq4\cr
3& e\in2\mathbf Z,\ e\geq4\cr
4& e\in2\mathbf Z,\ e\geq2\cr
6& e\in2\mathbf Z,\ e\geq2\cr
12& e\in2\mathbf Z,\ e\geq2.
\end{array}
\]
Over the particular field $\mathbf Q_3(\zeta_3)$, Theorem~\ref{thm:B-RD2}
already gives an example of defect six.
\end{theorem}
\begin{proof}
The special fibre is supersingular by Proposition~\ref{prop:B-RD3}.
Over an algebraically closed field of characteristic three it is
isomorphic to $y^2=x^3-x$. Its automorphisms are
\[
 (x,y)\longmapsto(u^2x+r,u^3y),\qquad
 u^4=1,\quad r^3-r=0.
\]
Thus its automorphism group has order twelve, is $C_3\rtimes C_4$,
and the only possible nontrivial subgroup orders are $2,3,4,6,12$.
The inertia image embeds in this group.

Over a good-reduction extension of ramification degree $m$, all eight
nonzero $3$-torsion parameters have positive integral valuations.
Their sum is $v_L(3)=me$, by the constant term of $[3](T)/T$.
Consequently $me\geq8$, which rules out $m=2,3$ at $e=2$.
The Weil pairing gives $\mu_3\subset K$, hence $e$ is even.
These are all the necessary restrictions in the table.

The tame cases $m=2,4$ are Theorem~\ref{thm:B-RD7}.
We explain the remaining constructions explicitly. For any desired
even $e$, a convenient base with algebraically closed residue field is
\[
 K_e=\operatorname{Frac}\left(
 W(\overline{\mathbf F}_3)[\varpi]/
 (\varpi^e+3\varpi^{e/2}+3)\right).
\]
The polynomial is Eisenstein, and $1+\varpi^{e/2}$ is a primitive
cube root of unity. In all the cases below the variable-ramification
family is taken over $K_e$.

For defect twelve and $e\geq4$, use $r=1$ unless $3\mid e+1$, in
which case use $r=3$. The second case first occurs at $e=8$, so always
$0<2r<e$. The integer $e+r$ is odd and prime to three. The missing
case $e=2$ is Theorem~\ref{thm:B-RD8}.

For defect six with $e\equiv0\pmod4$, use $r=2$, replacing it by
$r=6$ when $3\mid e+2$. These choices work for every such $e\geq8$.
For $e\equiv2\pmod4$, use $r=4$, replacing it by $r=8$ when
$3\mid e+4$. These choices work except at $e=2,6,14$.
In each applicable case, $e+r\equiv2\pmod4$ and $3\nmid e+r$,
so Lemma~\ref{lem:B-RD9} applies. The case $e=2$ is
Theorem~\ref{thm:B-RD2}. The case $e=4$ is obtained from the
order-twelve example by the tame quadratic extension: its inertia
image then has order six. The case $e=14$ follows from the order-six
example at $e=2$ by a tame degree-seven extension.
For $e=6$, start with the order-six example over a base of algebraically
closed residue field, and adjoin a cube root of a uniformizer.
Its good-reduction cubic subfield is obtained by adjoining a cube root
of $-1/4$, whose Kummer class has valuation zero; the uniformizer class
has valuation one. These cyclic cubic extensions are distinct and
linearly disjoint. Thus the order-six inertia image remains unchanged.
All finite extensions here are totally ramified.

For defect three with $e\equiv0\pmod4$, use $r=4$, replacing it by
$r=8$ when $3\mid e+4$. These choices work for every such $e\geq12$.
For $e\equiv2\pmod4$ and $e\geq6$, use $r=2$, replacing it by
$r=6$ when $3\mid e+2$; only $e=10$ is exceptional.
The examples at $e=4,8$ come from the order-six example at $e=2$
by tame extensions of degree two and four: they remove its order-two
inertia and retain the order-three subgroup.
For $e=10$, work over $K_{10}$ and put $t=1+\varpi^2$.
In the $2$-division cubic scaled by $x=3X$, translate
$X=\varpi^4+\varpi^6+Z$. Its constant, linear and quadratic
coefficients have valuations $22,16,14$, respectively. Its Newton
polygon therefore has the single slope $-22/3$, forcing a ramified
cubic extension. Here $v_K(\Delta)=108$, and the good-model
construction has $v_K(u^4)=36$; its fourth root requires no additional
ramification over the cubic extension. Hence the defect is exactly
three. The exact certificate \texttt{verify\_p3\_ranges.py} verifies these
coefficient valuations by reducing polynomials modulo
$\varpi^{10}+3\varpi^5+3$, and separately audits all the index recipes
for every even $e\leq10000$.

For the assertions about tame base change, one can work throughout
over completed maximal unramified bases. The order-six example has
cyclic inertia $C_6$. The order-twelve example has the full group
$C_3\rtimes C_4$, whose tame quotient is $C_4$; a tame quadratic
extension leaves the order-six subgroup. These observations also
verify that the stated defects are not merely upper bounds.
\end{proof}

\begin{corollary}[B-RD11: the cyclotomic field]\label{cor:B-RD11}
\status{Proved.}
For $p\geq5$, no elliptic curve over $\mathbf Q_p(\zeta_p)$ can have
full rational $p$-torsion and additive potentially good reduction.
For $p=3$, such a curve exists.
\end{corollary}
\begin{proof}
For $p\geq7$, the necessary inequality in
Theorem~\ref{thm:B-RD7} gives
$e\geq(p^2-1)/6>p-1$, contradicting $e(\mathbf Q_p(\zeta_p))=p-1$.
For $p=5$, the only possible defect is six. We justify that the good
extension may be chosen totally ramified, so its residue field really
is $\mathbf F_5$. Choose a minimal integral short Weierstrass equation
$y^2=x^3+ax+b$, and put $\delta=v_K(\Delta)$. Since the reduction is
potentially good, $v_K(j)\geq0$, giving $3v_K(a)\geq\delta$.
The discriminant identity then gives $2v_K(b)\geq\delta$ as well.
Put $m'=12/\gcd(12,\delta)$. Over the totally ramified extension
$K(\varpi^{1/m'})$, scaling by $u=\varpi^{\delta/12}$ gives integral
coefficients and unit discriminant. Conversely the discriminant
forces $m'$ to divide every good-reduction ramification degree.
Thus $m'=6$, and this construction gives the required degree-six
extension with unchanged residue field. Equality now holds in the
valuation bound: over this good-reduction extension every nonzero
$5$-torsion formal parameter has valuation one. Dividing those
parameters by a uniformizer and reducing gives an injective
$\mathbf F_5$-linear map
\[
 E[5](K)\longrightarrow k_L=k_K=\mathbf F_5.
\]
Additivity follows from the linear term of the formal group law;
injectivity follows because every nonzero parameter has valuation
one. Its source would have dimension two and its target dimension
one, a contradiction. The claim for $p=3$ is
Theorem~\ref{thm:B-RD2}.
\end{proof}

\subsection{All torsion levels for tame potentially good elliptic curves}

\begin{theorem}[B-RD12: the sharp tame higher-level criterion]\label{thm:B-RD12}
\status{Proved; independently audited by the parent agent.}
Fix $p\geq5$, $n\geq1$, and $m\in\{2,3,4,6\}$ with $m\mid p+1$.
Assume $e$ is a positive multiple of $p-1$, and put
$s=v_p(e/(p-1))$. There is a complete discretely valued $K$ of absolute
ramification index $e$ and an elliptic curve with full rational
$p^n$-torsion and semistability defect $m$ if and only if
\[
 s\geq n-1,\qquad
 e\geq\frac{p^2-1}{m}\,
        p^{\max(2n-2,\,4n-4-s)}.                  \tag{B.12.1}
\]
Arbitrary perfect residue fields are allowed; the constructions use
$\overline{\mathbf F}_p$. The same assertion holds for $p=3$ with a
specified tame defect $m=2$ or $4$. In that assertion no restriction on
wild defects is claimed.
\end{theorem}
\begin{proof}
We first obtain an exact integer criterion. Pass to completed maximal
unramified bases and a tame cyclic degree-$m$ good-reduction extension
$L/K$. Use the eigenparameter from Theorem~\ref{thm:B-RD7}.
Let $A$ and $B$ be the maximum and minimum valuations of a primitive
$p^n$-torsion point. Points $P,Q$ achieving them form a basis if $A>B$:
otherwise $Q=aP+pR$ with $a$ a unit; multiplication by $p$ strictly
increases every positive formal valuation, contradicting minimality
of $B$. If $A=B$, choose any basis. We have
\[
 A\geq B\geq1,\qquad A\equiv B\equiv q\pmod m,
 \qquad \gcd(q,m)=1.                              \tag{B.12.2}
\]

The full-level deformation ring has the factorization of
\cite[Proposition~2.2.3(vii)]{CarterStrauch2021LubinTate}.
Iterated $j$ times, it expresses the coordinate of $[p^j]Q$ as a unit
times a product of $p^j$ primitive level-$n$ coordinates, each raised
to $p^j$. The relevant lines have representatives $(a,1)$ with
$a\equiv0\pmod {p^{n-j}}$. All these coordinates have valuation $B$,
because the contribution $[a]P$ has valuation strictly greater than
$B$. Consequently
\[
 v_L(t([p^j]Q))=p^{2j}B\quad(0\leq j\leq n-1).   \tag{B.12.3}
\]
This factorization applies to the actual good elliptic curve by its
map from the universal framed deformation ring with full level.

Primitive vectors modulo scalar units are the lines of
$\mathbf P^1(\mathbf Z/p^n)$. There are $p^n$ representatives $(a,1)$,
all of valuation $B$: when $A=B$, a cancellation would contradict the
chosen maximum $A$. There is the line $(1,0)$, of valuation $A$.
For $1\leq j<n$, there are $(p-1)p^{n-j-1}$ representatives $(1,b)$
with $v_p(b)=j$, each of valuation $\min(A,p^{2j}B)$. Again equality
cannot cause cancellation beyond $A$, since $A$ is the maximum among
primitive points. The series $[p^n](T)/[p^{n-1}](T)$ has constant term
$p$, distinguished degree $(p^2-1)p^{2n-2}$, and precisely the
primitive $p^n$-torsion roots. Taking its constant term gives
\[
 me=(p-1)p^{n-1}\left(
 A+p^nB+\sum_{j=1}^{n-1}(p-1)p^{n-j-1}
                \min(A,p^{2j}B)\right).          \tag{B.12.4}
\]

We prove that every pair of integers satisfying (B.12.2) and (B.12.4)
can be realized; for the stated range it will suffice to use $B=q=1$.
Use the full Drinfeld $p^n$-level deformation ring $R_n$, whose regular
parameters $x,y$ are the universal basis coordinates
\cite[Theorem~2.1.2 and Proposition~2.2.3(i)]{CarterStrauch2021LubinTate}.
The eigenparameter averaging and equivariant Cohen presentation of
Theorem~\ref{thm:B-RD7} apply verbatim to this level.
We justify the weighted specialization, including its possible
cancellations, as follows.

Work modulo $p$ in the power-series presentation and give $X,Y$
weights $A,B$. Lift the full-level factorization (B.12.3), rather than
an arbitrary expression for $[p^j]y$. Each factor with representative
$(a,1)$ and $a\equiv0\pmod {p^{n-j}}$ has weighted initial form $Y$,
up to a nonzero constant. Therefore a lift of $[p^j]y$ has weighted
initial form $c_jY^{p^{2j}}$, with $c_j\in k^\times$. It follows that
for a primitive line $(1,b)$, $v_p(b)=j$, one can choose a lift whose
initial form is the lower-weight term of
\[
 X+c_bY^{p^{2j}},\qquad c_b\ne0,
\]
or that entire binomial when the two weights agree. The other
primitive lines are handled by their linear terms. These choices are
legitimate lifts because the factorization identities hold in $R_n$;
changes by the defining relation do not alter their images.

The constant-term identity for $[p^n](T)/[p^{n-1}](T)$ gives a generator
of the Cohen kernel of the form $p$ minus a unit times the product of
all primitive coordinates. Average this generator under $C_m$.
Its weighted initial form modulo $p$ is preserved: every initial
factor has scalar character $\zeta$, since
$p^{2j}\equiv1\pmod m$, and the total number
$(p^2-1)p^{2n-2}$ of factors is divisible by $m$.
Now specialize $X=\pi^A$, $Y=c\pi^B$, choosing the residue of $c$
to avoid the finitely many roots of the displayed initial forms.
The resulting relation has linear term $p$ and first nonzero term
modulo $p$ of degree $d$, where $d$ is the right side of (B.12.4).
It defines a complete DVR with $v_L(p)=d$.
Every primitive-coordinate weight is strictly smaller than $d$, so
terms in its chosen lift divisible by $p$ cannot change that weight.
Thus all the asserted valuations hold after specialization, with no
circular assumption on $v_L(p)$.

For $A\equiv B\equiv1\pmod m$, the action $\pi\mapsto\zeta\pi$
preserves the relation and the specialization. Tame descent gives a
curve over the invariant field, of absolute ramification $d/m$,
with full rational $p^n$-torsion and exact defect $m$, as in
Theorem~\ref{thm:B-RD7}. For a general common residue $q$ in
(B.12.2), replace the action on $\pi$ by multiplication by
$\zeta^{q^{-1}}$, where the inverse is taken modulo $m$.

It remains to simplify (B.12.4), which also proves that $B=1$ suffices.
Choose $r\in\{0,\ldots,n-1\}$ so that
$p^{2r}B\leq A\leq p^{2r+2}B$, with the upper bound omitted for
$r=n-1$. Summing the geometric series in (B.12.4) gives
\[
 me=(p-1)p^{2n-r-2}(A+p^{2r+1}B).                \tag{B.12.5}
\]
It implies the two conditions
\[
 p^{2n-r-2}(p-1)\mid e,\qquad
 e\geq \frac{p^2-1}{m}p^{2n+r-2}.                \tag{B.12.6}
\]
Conversely suppose these conditions hold for some $r$. Increase $r$
until it is $n-1$, or until the second inequality fails for $r+1$;
the divisibility remains true. Put
\[
 B=1,\qquad
 A=\frac{me}{(p-1)p^{2n-r-2}}-p^{2r+1}.
\]
Then $A$ is an integer, $A\equiv1\pmod m$, and
$p^{2r}\leq A<p^{2r+2}$ unless $r=n-1$, when only the lower bound
is needed. Hence it realizes (B.12.4) by the construction above.
Finally the least permissible $r$ is
$\max(0,2n-2-s)$; such an $r\leq n-1$ exists exactly when
$s\geq n-1$. Substitution gives (B.12.1).
\end{proof}

\begin{corollary}[B-RD13: the least good-forcing level in the tame case]
\label{cor:B-RD13}
\status{Proved; independently audited by the parent agent.}
Let $p\geq5$, $e\in(p-1)\mathbf Z_{>0}$, and let $m_{\max}(p)$ be
as in Theorem~\ref{thm:B-RD7}. If $m_{\max}(p)e<p^2-1$, put $N=0$.
Otherwise put
\[
 s=v_p(e/(p-1)),\qquad
 L=\left\lfloor\log_p\frac{m_{\max}(p)e}{p^2-1}\right\rfloor,
 \qquad
 N=1+\min\left(s,\left\lfloor\frac L2\right\rfloor,
                     \left\lfloor\frac{L+s}{4}\right\rfloor\right).
\]
Among potentially good elliptic curves over all such fields of
absolute ramification $e$, the least positive level $n$ at which
full rational $p^n$-torsion forces good reduction is $N+1$.
For $p=3$, the same formula with $m_{\max}=4$ gives the least level
within the class of potentially good curves with tame inertia image.
\end{corollary}
\begin{proof}
For a fixed level, the largest allowed defect gives the weakest
inequality in Theorem~\ref{thm:B-RD12}. Its inequalities are equivalent
to $n-1\leq s$, $2n-2\leq L$, and $4n-4-s\leq L$.
Solving them gives exactly the displayed largest bad level $N$.
The construction realizes it, and its lower levels, while every
potentially good curve outside the bad cases is good.
\end{proof}

\begin{remark}[The scope of ``forces good reduction'']
The potentially good hypothesis is essential. If $K$ contains
$\mu_{p^n}$, a split Tate curve with parameter $q=a^{p^n}$,
$0<|a|<1$, has full rational $p^n$-torsion and split multiplicative
reduction. Thus the preceding least-level assertion is not a
criterion for all elliptic curves without that hypothesis.
\end{remark}

\subsection{Cyclotomic bounds in every dimension and a character lower bound}

\begin{lemma}[B-RD14: full torsion forces the cyclotomic extension]
\label{lem:B-RD14}
\status{Proved.}
Let $A/K$ be an abelian variety of positive dimension over a
characteristic-zero field. If $A[p^n]$ is rational over $K$, then
$\mu_{p^n}\subset K$. A principal polarization, or a polarization of
degree prime to $p$, is not needed for this implication.
\end{lemma}
\begin{proof}
Choose any polarization defined over $K$. Its Weil pairing gives a
nonzero $G_K$-equivariant integral alternating form
\[
 b:T_pA\times T_pA\longrightarrow\mathbf Z_p(1).
\]
Divide all its coefficients by their largest common power of $p$.
The resulting form is still integral and equivariant, and at least
one coefficient is a unit. Since $G_K$ acts trivially on $T_pA/p^n$,
equivariance yields
$(\chi_p(\sigma)-1)b(x,y)\equiv0\pmod {p^n}$ for every $x,y$.
Choose a pair with unit value; then
$\chi_p(\sigma)\equiv1\pmod {p^n}$ for every $\sigma\in G_K$.
This is precisely the claimed cyclotomic containment.
\end{proof}

\begin{proposition}[B-RD15: the unconditional reading of the least-level question]
\label{prop:B-RD15}
\status{Proved; the final implication is explicitly vacuous.}
Fix $p\geq3$, a positive multiple $e$ of $p-1$, and $g\geq1$, and put
$s=v_p(e/(p-1))$. Among \emph{all} abelian varieties of dimension $g$
over all complete discretely valued fields with perfect residue field
and absolute ramification $e$, the least positive $n$ for which full
rational $p^n$-torsion implies good reduction is $s+2$.
At this level full rational $p^n$-torsion cannot occur at all.
Consequently a substantive good-reduction threshold must specify the
potentially good class, as in Corollary~\ref{cor:B-RD13}.
\end{proposition}
\begin{proof}
Lemma~\ref{lem:B-RD14} and the ramification degree
$[\mathbf Q_p(\mu_{p^n}):\mathbf Q_p]=p^{n-1}(p-1)$ imply
$p^{n-1}(p-1)\mid e$. Thus full rational $p^n$-torsion is impossible
for $n\geq s+2$.

Conversely, for each $1\leq n\leq s+1$, choose a complete field $K$
with residue field $\overline{\mathbf F}_p$, absolute ramification $e$,
and $\mu_{p^n}\subset K$. Such a field is obtained by taking the
cyclotomic extension of $\operatorname{Frac}W(\overline{\mathbf F}_p)$
and a further totally ramified extension of the required degree.
By \cite[Theorem~1, p.6 of the author version]{Tate1995Uniformization},
a split Tate curve with parameter $q=a^{p^n}$, where $a$ is a
uniformizer, has all its $p^n$-torsion represented by
$\zeta_{p^n}^i a^j$ in $K^\times/q^{\mathbf Z}$, and has split
multiplicative reduction. For dimension $g$, take its product with
$g-1$ canonical ordinary elliptic lifts. Their $p$-divisible groups
are $\mu_{p^\infty}\oplus\mathbf Q_p/\mathbf Z_p$, so they too have
full rational $p^n$-torsion. The product still has bad reduction.
This proves the asserted least level and explains its vacuity.
\end{proof}

\begin{proposition}[B-RD16: necessary lower bounds for character determination]
\label{prop:B-RD16}
\status{Proved.}
Let $h\geq2$. For $p\geq5$ and absolute ramification $e$, let $N$
be the largest bad elliptic level from Corollary~\ref{cor:B-RD13}.
If $N\geq1$, any universal least index $\nu$ in B.4 must satisfy
$\nu\geq N$. More precisely, for each $1\leq n\leq N$ there are two
height-$h$ BT groups satisfying the manuscript's monodromy condition,
with isomorphic $p^n$-torsion, whose characters on inertia elements of
order prime to $p$ differ. For $p=3$, the same lower bound holds using
the tame value $m_{\max}=4$ in that corollary.
\end{proposition}
\begin{proof}
Take the elliptic curve $E/K$ supplied by Theorem~\ref{thm:B-RD12}
at level $n$, and a canonical ordinary elliptic lift $E'/K$.
Lemma~\ref{lem:B-RD14} gives $\mu_{p^n}\subset K$, so both
$E[p^n]$ and $E'[p^n]$ are constant of rank $p^{2n}$.
Their inertia characters differ: the ordinary good curve has trace
two, whereas the nontrivial tame automorphism of the supersingular
good reduction has trace different from two. Append the same constant
etale BT group of height $h-2$ to both. Their characters each increase
by $h-2$, preserving the difference, and their truncated groups remain
isomorphic. These groups are potentially good, hence satisfy the
stated monodromy condition without using any extension argument.
Finally the manuscript's $M(\nu)$ is $M[p^{\nu+1}]$; failure at
$n=N$ gives $\nu\geq N$.
\end{proof}

\subsection{The monodromy condition and the Fontaine comparison}

\begin{proposition}[B-RD4: Lead B' needs an additional hypothesis]\label{prop:B-RD4}
\status{Refuted as stated.}
For every odd prime $p$, there is a height-two Barsotti--Tate group over
$K=\mathbf Q_p(\zeta_p)$ satisfying the manuscript's monodromy condition
whose rational Tate module is not potentially Hodge--Tate, and therefore
is not potentially semistable. Its manuscript inertia character is the
constant function $2$. Thus the monodromy condition in 9:11 cannot be
identified with potential semistability, and Fontaine's usual
Weil--Deligne representation need not be defined for its Tate module.
\end{proposition}
\begin{proof}
Write $\chi:G_K\to\mathbf Z_p^\times$ for the cyclotomic character.
Since $\mu_p\subset K$, it takes its values in $1+p\mathbf Z_p$. Put
\[
 \rho(g)=\begin{pmatrix}1&p^{-1}\log\chi(g)\\0&1\end{pmatrix}.
\]
This is a continuous representation on $T=\mathbf Z_p^2$, because the
logarithm is additive on $1+p\mathbf Z_p$ and takes values in
$p\mathbf Z_p$. Over a characteristic-zero field, the equivalence
between finite etale group schemes and finite continuous Galois modules
applied to $T/p^nT$ gives compatible finite etale group schemes $M[p^n]$.
The transition maps are induced by multiplication by $p$ on the lattices,
and the resulting inductive system is a height-two Barsotti--Tate group.
The invariant first coordinate and the quotient both give the constant
group $\mathbf Q_p/\mathbf Z_p$. They extend over $\mathcal O_K$, so the
specified two-step filtration satisfies the manuscript condition and
its character is $1+1=2$.

For $V=T\otimes\mathbf Q_p$, the Sen operator is
\[
 \Theta_V=\frac1p\begin{pmatrix}0&1\\0&0\end{pmatrix}.
\]
Indeed the logarithm of a unipotent matrix $1+aN$, with $N^2=0$, is
$aN$, and division by $\log\chi(g)$ gives the displayed operator.
It remains nonzero and nilpotent after restriction to any finite
extension. Sen's criterion says that a representation is Hodge--Tate
exactly when this operator is diagonalizable with integral eigenvalues;
the identical unscaled example and the criterion are explicitly given
in \cite[\S II.1.2, p.~9 of the author's version]{Berger2004Introduction}.
The operator is not diagonalizable. Potentially semistable
representations are potentially Hodge--Tate, so no finite extension
makes this representation semistable.
\end{proof}

\begin{remark}[The corrected scope of the proposed dictionary]
If one separately assumes potential semistability of $T_p(M)\otimes
\mathbf Q_p$, the comparison proposed in Lead B' is the appropriate
one: on a potentially good graded piece, crystalline comparison
identifies its Dieudonne isocrystal with the corresponding period
module, equivariantly for descent. Exactness of the potentially
semistable period functor then adds these characters across a
filtration. Covariant and contravariant conventions must be fixed:
with $D_{\mathrm{cov}}(G)=D_{\mathrm{contra}}(G^D)$, the covariant
isocrystal corresponds to $D_{\mathrm{cris}}(T_pG)(-1)$; the Tate twist
has trivial Weil--Deligne inertia character. This paragraph records
the corrected comparison to be sourced; the counterexample above
requires no comparison theorem.
\end{remark}

\begin{proposition}[B-RD5: the height-one case of B.4]\label{prop:B-RD5}
\status{Proved.}
Suppose $p\geq3$ and $\mu_p\subset K$. For height $h=1$, the manuscript
character on prime-to-$p$ inertia elements is determined by $M(0)$.
The least index in its indexing is therefore $\nu=0$.
\end{proposition}
\begin{proof}
A height-one group satisfying the monodromy condition itself acquires
good reduction. Its good model has dimension zero or one. In dimension
zero it is etale, and in dimension one its Cartier dual is etale.
Over the completed maximal unramified extension of a good-reduction
field its model is respectively $\mathbf Q_p/\mathbf Z_p$ or
$\mu_{p^\infty}$. Its finite descent inertia character therefore takes
values in the torsion subgroup of $\mathbf Z_p^\times$, namely
$\mu_{p-1}$. In dimension zero, reduction modulo $p$ gives precisely
this character. In dimension one, the Tate module acquires additionally
the cyclotomic factor, whose reduction is trivial under $\mu_p\subset K$.
Thus in both cases the finite character is the unique Teichmuller lift
of the one-dimensional mod-$p$ inertia character. The same conclusion
holds after summing characters of height-one potentially good graded
pieces: the trace on a prime-to-$p$ element is the Brauer character of
the mod-$p$ representation, and is independent of the chosen filtration.
\end{proof}

